\section{Sua parte: análise exploratória}
\textbf{Mensagem principal:} a taxa é assimétrica à direita (o que pede uma transformação), varia muito entre regiões e estados, e nenhum fator sozinho a explica; renda, \% pobres e alfabetização medem quase o mesmo, então só a renda entra no modelo.

\subsection{O que enfatizar}
\begin{itemize}
\item Daqui em diante, tudo usa os \nMunEsc{} municípios do escopo (no mapa, os de fora aparecem em cinza).
\item \textbf{Resumo:} as variáveis têm escalas e dispersões bem diferentes; a densidade é muito assimétrica e por isso entra em log.
\item \textbf{Assimetria (histograma da taxa):} metade dos municípios fica entre \nTaxaQuartilUm{} e \nTaxaQuartilTres, mas há uma cauda longa à direita (assimetria \nAssimTaxa, positiva), que puxa a média (\nTaxaMedia) para cima da mediana (\nTaxaMediana). Isso antecipa a transformação do modelo.
\item \textbf{Variação entre regiões (mapa e boxplot no mesmo slide):} Norte e Nordeste acima de Sudeste e Sul. Daí o estado no modelo; a Amazônia Legal sai por ficar redundante com ele.
\item \textbf{Correlações:} a maior com o log da taxa é a de \nCorMaxVar{} (\nCorMax). Na matriz, renda, \% pobres e alfabetização se repetem (até \nCorTrio{} entre si): multicolinearidade, VIF da renda \nVIFrenda.
\end{itemize}

\subsection{Perguntas prováveis}
\begin{description}[leftmargin=0pt]
\item[Correlação alta com \% preta ou parda quer dizer que a cor causa violência?] Não. É associação: reflete desigualdade racial e a concentração da vitimização na população negra, como mostram o Atlas da Violência e a literatura.
\item[Por que o log e não a raiz?] Quem decide é o Box-Cox, no slide do Luiz Felipe; aqui só mostramos que a taxa original é assimétrica.
\item[Tirar \% pobres e alfabetização não perde informação?] Pouca: elas repetem a renda; mantê-las deixaria os coeficientes instáveis e com sinais trocados.
\item[A fronteira não importa?] No modelo ela não é significativa: com o estado no modelo, boa parte do efeito regional já é captada; a seleção de variáveis vai decidir se ela fica.
\end{description}

\subsection{Transição}
``A relação com a renda tem uma forma especial, que o Luiz Felipe mostra agora.''
